% Native semantic reconstruction of source physical page 13 (printed page 12).
% Provenance: IEEE Std 1516-2025, Clauses 1.4--1.4.3 and footnotes 8--9.
\HLAStartSemanticPage{13}{12}
\fontsize{10.02pt}{11.523pt}\selectfont
\begingroup
\enlargethispage{-11.8pt}
\hbadness=10000\relax
\addfontfeatures{FakeStretch=1.035}
\vspace*{-0.81pt}

\HLASubclauseHeading{clause-1.4}{1.4 Background}
\vspace{10.63pt}

The HLA is a technical architecture developed to facilitate the reuse and interoperation of simulation\linebreak
systems and assets (e.g., federation managers, data collectors, passive viewers, real world (“live”) systems,\linebreak
and other utilities). It was developed in response to the clear recognition of the value of modeling and\linebreak
simulation (M\&S) in supporting a wide variety of applications as well as the need to manage M\&S assets\linebreak
as cost effective tools. The selected definition of an architecture as intended in HLA---“major functional\linebreak
elements, interfaces, and design rules, pertaining as feasible to all simulation applications, and providing a\linebreak
common framework within which specific system architectures can be defined”---is one that is commonly\linebreak
accepted and is consistent with the definition of architecture for computer simulations provided by the\linebreak
IEEE.\par
\vspace{20.48pt}

\HLASubclauseHeading[10.02pt]{clause-1.4.1}{1.4.1 Components of the HLA standard}
\vspace{10.18pt}

The HLA provides a general framework within which developers can structure and describe their\linebreak
simulation systems and/or assets and interoperate with other simulation systems and assets. The HLA\linebreak
consists of three main components. The first is the HLA Framework and Rules (i.e., this standard), which\linebreak
provides a set of ten rules, five applying to \textit{federates} (see 3.1) and five applying to \textit{federations} (see 3.1),\linebreak
that together ensure the proper interaction of federates in a federation and define the responsibilities of\linebreak
federates and federations. The second component is a specification of services that allow simulations to\linebreak
connect to one another, exchange data, and coordinate activities during a distributed runtime execution.\linebreak
This component is the HLA federate interface specification (found in IEEE Std 1516.1\HLAStandardMark-2025), which\linebreak
describes the capabilities and principles of operation of an underlying software infrastructure, called the\linebreak
\textit{runtime infrastructure (RTI)}, that will be required for runtime operation. The third component is the HLA\linebreak
object model template (OMT) Specification (found in IEEE Std 1516.2\HLAStandardMark-2025).\textsuperscript{8} It provides a\linebreak
specification for describing \textit{object models} (see 3.1) that define the information produced or required by a\linebreak
simulation application and for reconciling definitions among simulations to produce a common data model\linebreak
for mutual interoperation.\par
\vspace{20.51pt}

\HLASubclauseHeading[10.02pt]{clause-1.4.2}{1.4.2 Terminology}
\vspace{10.19pt}

Some terminology is necessary to give an overview of the HLA framework. A federate application is an\linebreak
application that supports the HLA interface to a RTI and that is capable of joining a federation execution.\linebreak
Note that, in addition to a simulation, a federate application may represent an interface to a live system or a\linebreak
support utility such as an event logger or performance monitor. A federate application may be coupled with\linebreak
other software applications and a RTI. A joined federate is a member of a federation execution, actualized\linebreak
by a federate application invoking the \HLAService{Join Federation Execution} service. A federation is a named set of\linebreak
federate applications and a common federation object model (FOM) that are used as a whole to achieve\linebreak
some specific objective. A federation execution is the actual operation, over time, of a set of joined\linebreak
federates that are interconnected by an RTI.\par
\vspace{12.23pt}

NOTE---Formal definitions for each of these terms can be found in the definitions clause of all three HLA\linebreak
specifications.\textsuperscript{9}\par
\vspace{18.29pt}

\HLASubclauseHeading[10.02pt]{clause-1.4.3}{1.4.3 Object models}
\vspace{10.18pt}

There are three types of object models in the HLA framework. First is a simulation object model (SOM).\linebreak
The SOM is a specification of the types of information that an individual federate application could provide\linebreak
to HLA federations as well as the information that an individual joined federate can receive from other\linebreak
joined federates in HLA federation executions. The SOM may be viewed as a federate application’s native\par

\footnotetext[8]{Information on references can be found in Clause 2.}
\footnotetext[9]{Notes in text, tables, and figures are given for information only and do not contain requirements needed to implement the standard.}
\endgroup
\clearpage
